
%%%%%%%%%%%%%%%%%%%%%%%%%%%%%%%%%%%%%%%%%%%%%%%%%%%%%%%%%%%%%
\section{符号说明}

本文使用的主要符号及其含义如表\ref{tab:符号说明}所示，未列出的符号在首次出现处给出说明。

\begin{table}[H]
	\centering
	\caption{符号说明}
	\begin{tabularx}{\textwidth}{CL}
		\toprule
		符号    & 说明    \\
		\midrule
		$R,\ D$ & 目标区域半径（1800 米）与目标圆域 \\
		$\delta$ & 示向度误差上界，$\delta=1^\circ=0.0174533$ 弧度 \\
		$S_i,\ \theta_i$ & 第 $i$ 个检测点及其测得的示向度 \\
		$W_i$ & 由 $S_i$ 及 $\theta_i$ 确定的张角 $2\delta$ 的楔形区域 \\
		$P$ & 交会定位区域（各楔形与目标圆域的交） \\
		$D_P$ & 定位区域直径，$D_P=\max_{a,b\in P}\lVert a-b\rVert$ \\
		$G,\ G_c$ & 干扰源位置与频道 $c$ 对应的干扰源位置 \\
		$r_c,\ R_c^{\min}$ & 干扰源的有效接收半径与其下界（1000 米） \\
		$\Omega$ & 由单次示向度确定的干扰源可行域 \\
		$b,\ \varphi$ & 基线长度与基线方位角与视线方向的夹角 \\
		$r_1,\ r_2,\ \alpha$ & 两检测点到干扰源的距离与干扰源处的交会角 \\
		$E(S_2;G)$ & 由两站方位线扰动引起的定位误差上界 \\
		$\mathcal C$ & 第二检测点的候选区域 \\
		$\mathcal Q,\ N_q$ & 观测点集及其元素个数 \\
		$d,\ \rho(d)$ & 覆盖网的环半径与其对应的覆盖半径 \\
		$L_{\mathrm{TSP}}$ & 清除阶段的旅行商巡回长度 \\
		$T$ & 任务完成的虚拟总时间 \\
		$N_m,\ N_{sw}$ & 检测次数与频道切换次数 \\
		$\psi_c$ & 定向源 $c$ 的定向方向 \\
		$H_c$ & 定向源的可检测域（半圆盘） \\
		$\kappa$ & 定向场景相对全向场景的时间惩罚系数 \\
		\bottomrule
	\end{tabularx}

	\label{tab:符号说明}
\end{table}
